\documentclass[10pt,a4paper]{article}

\usepackage[a4paper,top=1.15cm,bottom=1.15cm,left=1.25cm,right=1.25cm]{geometry}
\usepackage{fontspec}
\usepackage[english]{babel}
\usepackage{hyperref}
\usepackage{xcolor}
\usepackage{titlesec}
\usepackage{enumitem}

\setmainfont{Arial}
\pagestyle{empty}
\setlength{\parindent}{0pt}
\setlength{\parskip}{2pt}
\linespread{1.02}
\hyphenpenalty=10000
\exhyphenpenalty=10000
\sloppy
\setlist[itemize]{leftmargin=1.2em,itemsep=2.2pt,topsep=2.5pt,parsep=0pt,label=-}
\urlstyle{same}

\definecolor{heading}{HTML}{123867}
\definecolor{linkblue}{HTML}{005EB8}
\hypersetup{colorlinks=true,urlcolor=linkblue,linkcolor=linkblue}
\titleformat{\section}{\large\bfseries\color{heading}}{}{0pt}{}[\vspace{-4pt}\rule{\linewidth}{0.45pt}]
\titlespacing{\section}{0pt}{8pt}{5pt}

\newcommand{\entry}[3]{\textbf{#1} \hfill #3\\\textit{#2}\vspace{1pt}}
\newcommand{\project}[3]{\textbf{#1} \hfill #3\\\textit{#2}\vspace{1pt}}

\begin{document}
\fontsize{10}{11.7}\selectfont

\begin{minipage}[t]{0.60\textwidth}
{\LARGE \textbf{Kossi Richard Allado}}\\[2pt]
\textbf{Engineering Student in Cybersecurity and Artificial Intelligence}\\[1pt]
Cybersecurity \textbar{} SOC/DFIR \textbar{} Offensive Security \textbar{} AI Security
\end{minipage}
\hfill
\begin{minipage}[t]{0.36\textwidth}
\raggedleft
Casablanca, Morocco\\
+212 621 074 305\\
alladokossirichard2025@gmail.com
\end{minipage}

\vspace{4pt}
\footnotesize
\href{https://github.com/ALLAKORI}{GitHub} \quad | \quad
\href{https://allakori.github.io/}{Portfolio} \quad | \quad
\href{https://www.linkedin.com/in/kossi-richard-allado-50177b2a5}{LinkedIn}
\normalsize

\section{Profile}
Engineering student in Cybersecurity and Artificial Intelligence at ENSA Beni Mellal, preparing for a PFE 2027 internship focused on cybersecurity. I am interested in SOC/DFIR, threat intelligence, application security, IAM, Linux security, Offensive Security labs, and AI-assisted security analysis. Earliest available date for PFE: February 1, 2027.

My profile is cyber-first: incident investigation, evidence collection, controlled exploitation, access-control reasoning, remediation documentation, and security-oriented backend/RAG work. AI is presented as a support capability for cyber analysis, documentation, and secure automation, not as a separate data-science target.

\section{Technical Skills}
\textbf{Cybersecurity research:} threat intelligence monitoring, alert triage, IOC extraction, incident timelines, MITRE ATT\&CK mapping -- Splunk, Sysmon, Event Viewer, Linux logs.\\
\textbf{SOC and DFIR:} log correlation, timeline reconstruction, evidence collection, containment and remediation reporting -- Splunk, Sysmon, Event Viewer, auditd, Linux logs.\\
\textbf{Application, IAM and platform security:} authentication review, RBAC, MFA/TOTP, rate limiting, account lockout, audit logs, secure configuration review -- JWT, Keycloak, Docker Compose, Linux, OpenSSL.\\
\textbf{Offensive Security validation:} reconnaissance, controlled exploitation, impact analysis and remediation documentation in lab environments -- Nmap, Burp Suite, Hydra, Netcat, Kali Linux.\\
\textbf{AI Security and secure automation:} AI-assisted threat analysis, RAG security controls, structured documentation and response guidance -- FastEmbed, Qdrant, LLM APIs, Python, SQL.

\section{Experience}
\entry{PFA Cybersecurity Intern -- Early Warning System for Cyber Threats Targeting SMEs}{CMRPI - Espace Maroc Cyberconfiance}{July -- August 2026}
\begin{itemize}
  \item Developed the \textbf{CyberVeille PME} MVP: a Streamlit application that lets SME users analyze suspicious URLs without opening them, with verdict, risk score, reasons, and recommendations.
  \item Integrated a detection chain combining local URLHaus threat-intelligence lookup, 21 URL features, scikit-learn Logistic Regression, SQLite history, and configurable SMTP alerts.
  \item Documented limitations and SME reflex sheets; validated the MVP with automated tests covering invalid inputs, URLHaus, CSV processing, SQLite, disabled alerts, and SMTP failure.
\end{itemize}

\entry{Artificial Intelligence / Predictive Maintenance Intern}{ONEE - Branche Eau, Kasba Tadla}{August 2025}
\begin{itemize}
  \item Prepared and analyzed high-pressure pump sensor data to identify potential fault signals.
  \item Trained a Random Forest model with Python, Pandas, and Scikit-learn to classify fault risk.
  \item Documented model evaluation with classification report, confusion matrix, feature importance, ROC, and precision-recall curves.
  \item Built a desktop test interface and documented the predictive-maintenance prototype.
\end{itemize}

\section{Projects}
\project{Chatbot USMS v2.0.0}{Academic product, ENSA Beni Mellal / Sultan Moulay Slimane University}{Secure RAG / Backend / Access Control}
\href{https://github.com/ALLAKORI}{GitHub}
\begin{itemize}
  \item Built a multilingual academic assistant providing natural-language access to university FAQs, timetables, and administrative information in French, English, and Arabic.
  \item Designed a RAG pipeline over PDF/JSON/TXT/MD sources using FastEmbed embeddings, Qdrant vector search, retrieval confidence, and LLaMA-3.3-70B generation through GroqCloud.
  \item Integrated FastAPI, Next.js 15, PostgreSQL, Redis, Docker Compose, SSE streaming, JWT/RBAC, rate limiting, account lockout, PII masking, audit logs, and automated tests.
\end{itemize}

\section{Cybersecurity Projects}
\project{Linux Server Compromise \& Forensic Investigation Lab}{Kali Linux attacker / Ubuntu target, SSH, UFW, auditd, Fail2ban}{Offensive Security / DFIR}
\href{https://github.com/ALLAKORI/Linux-Server-Compromise-Forensic-Investigation-Lab}{GitHub}
\begin{itemize}
  \item Built and documented a controlled scenario: reconnaissance, SSH brute force, initial access, malicious user creation, GTFOBins privilege escalation, persistence, and reverse shell.
  \item Collected defensive evidence: SSH connections, executed commands, sudo activity, created accounts, network flows, auditd artifacts, reverse shell port, and attack timeline.
  \item Validated remediation steps: UFW blocking, account removal, password reset, sudo/SSH review, and Fail2ban activation.
\end{itemize}

\project{Volt Typhoon Intrusion Investigation}{TryHackMe, Splunk, simulated Windows and application logs}{SOC / Threat Intelligence}
\href{https://allakori.github.io/writeups/tryhackme/volt-typhoon/}{Write-up}
\begin{itemize}
  \item Investigated a simulated intrusion involving compromised account activity, PowerShell/WMIC, web shell usage, Mimikatz, C2 communication, and log clearing.
  \item Reconstructed incident phases, extracted IOCs, built a timeline, and mapped observed activity to MITRE ATT\&CK tactics and techniques.
  \item Proposed detection and remediation ideas to improve monitoring and incident response.
\end{itemize}

\section{Education}
\textbf{Engineering Cycle -- Cybersecurity and Artificial Intelligence} \hfill 2024 -- Present\\
Sultan Moulay Slimane University -- National School of Applied Sciences, Beni Mellal\\
\textbf{DEUST -- Mathematics, Computer Science and Physics, with honors} \hfill 2022 -- 2024\\
Sidi Mohamed Ben Abdellah University -- Faculty of Sciences and Technologies, Fes\\
\textbf{Scientific Baccalaureate, Series D, high honors} \hfill 2022\\
Lycee de Djagble, Togo

\section{Certifications and Achievements}
\textbf{Certifications:} Google Cybersecurity Professional Certificate, Cisco Ethical Hacker, TryHackMe AI Security, TryHackMe Junior Penetration Tester, IBM Introduction to Data Engineering.\\
\textbf{Achievements:} TryHackMe Top 1\%, CSIA CTF 2026: 1st place with Team GHOSTSHELL, CiteFlag 2026: finalist, 6th place with Team CHILA, MACC 2026: participant.

\section{Languages}
French: fluent \textbar{} English: professional technical level \textbar{} Ewe: native

\end{document}
